\chapter{Research Methodology}\label{ch:method}

\section{Research approach}
The study follows a \emph{design-science} and \emph{quantitative} approach. Its main output is an artefact, the
analytics framework, which is built to solve a stated business problem and is then evaluated rigorously against
explicit criteria. This matches the guidelines of \citet{hevner2004} for design-science research in information
systems: the artefact must address a relevant problem, be evaluated rigorously, and be communicated to both
technical and managerial audiences. The empirical analysis is quantitative and uses a large secondary dataset of sales
transactions. The overall lifecycle follows CRISP-DM \citep{shearer2000} as refined into ten phases in the
project proposal (\cref{fig:lifecycle}).

\begin{figure}[htbp]\centering
\begin{tikzpicture}[node distance=4mm and 4mm, every node/.style={font=\small},
  box/.style={draw=brand, rounded corners, fill=brand!8, minimum width=3.3cm, minimum height=9mm, align=center},
  arr/.style={-{Stealth[length=2mm]}, thick, brand}]
\node[box] (p1) {1. Problem\\definition};
\node[box, right=of p1] (p2) {2. Data\\collection};
\node[box, right=of p2] (p3) {3. Data\\preparation};
\node[box, right=of p3] (p4) {4. Exploratory\\analysis};
\node[box, below=9mm of p4] (p5) {5. Feature\\engineering};
\node[box, left=of p5] (p6) {6. Model\\development};
\node[box, left=of p6] (p7) {7. Model\\evaluation};
\node[box, left=of p7] (p8) {8. Dashboard\\development};
\node[box, below=9mm of p8] (p9) {9. System\\evaluation};
\node[box, right=of p9] (p10) {10. Documentation};
\draw[arr] (p1)--(p2); \draw[arr] (p2)--(p3); \draw[arr] (p3)--(p4); \draw[arr] (p4)--(p5);
\draw[arr] (p5)--(p6); \draw[arr] (p6)--(p7); \draw[arr] (p7)--(p8); \draw[arr] (p8)--(p9); \draw[arr] (p9)--(p10);
\draw[arr, dashed, accent] (p7.north) to[bend left=20] node[above, font=\scriptsize]{iterate} (p3.south);
\end{tikzpicture}
\caption{Project lifecycle, adapted from CRISP-DM}\label{fig:lifecycle}
\end{figure}

\section{Problem definition}
Four analytical problems were derived from the business questions:
\begin{enumerate}
  \item \textbf{Demand forecasting (regression).} Predict the number of units of each product sold in the
        next month. This supports production and inventory planning.
  \item \textbf{Large-order prediction (classification).} At the moment an order is entered, predict whether
        it will be large (\(\geq 10\) units). This supports allocation of stock and delivery capacity.
  \item \textbf{Segmentation (clustering).} Group products and market cells (region $\times$ channel $\times$
        customer type) with similar behaviour, to support portfolio and sales-force strategy.
  \item \textbf{Anomaly detection.} Detect unusual trading days and suspicious transactions. This supports
        monitoring and data-quality control.
\end{enumerate}

\section{Data}
The dataset comprises \NRows transactions with \NCols attributes, recorded between \DateMin and \DateMax. Each
record holds an identifier and date; the product (identifier, name, category, pack size, launch year); price
information (printed maximum retail price, list price, market price, discount percentage and unit selling
price); the quantity; the gross, discount and net sales values in Bangladeshi Taka (BDT); and the region,
channel and customer type. \Cref{tab:dictionary} in \cref{app:dictionary} gives the full data dictionary. The
data contain no personal identifiers.

\section{Evaluation design}
\subsection{Temporal hold-out}
Because the problems involve time, every supervised model is trained on the years \textbf{2020--2024} and
tested on \textbf{January 2025 -- September 2026}. This unseen period of \TestMonths months plays the role of
``the future''. Random train/test splits were deliberately not used, because they would let the models learn
from observations that come after the ones they predict \citep{bergmeir2012}.

\subsection{Hyper-parameter tuning}
Hyper-parameters were tuned only on the training years, using an \emph{expanding-window} validation. For
validation year $y \in \{2022, 2023, 2024\}$, the model is fitted on all months before $y$ and scored on $y$.
The configuration with the best mean validation score is then refitted on 2020--2024 and evaluated once on the
test period. For classification, the probability threshold was chosen in the same way: the model was fitted on
2020--2023 and the threshold that maximised F1 on 2024 was kept.

\subsection{Metrics}
For regression, the following measures were used, where $y_i$ is the actual and $\hat y_i$ the predicted value
for $n$ product-months:
\begin{align*}
\text{MAE} &= \tfrac{1}{n}\textstyle\sum_i |y_i-\hat y_i|, &
\text{RMSE} &= \sqrt{\tfrac{1}{n}\textstyle\sum_i (y_i-\hat y_i)^2}, &
R^2 &= 1-\frac{\sum_i (y_i-\hat y_i)^2}{\sum_i (y_i-\bar y)^2},\\
\text{MAPE} &= \tfrac{100}{n}\textstyle\sum_i \frac{|y_i-\hat y_i|}{\max(y_i,1)}, &
\text{WAPE} &= 100\,\frac{\sum_i |y_i-\hat y_i|}{\sum_i y_i}. &&
\end{align*}
WAPE is preferred for business communication because it weights each product by its volume. For
classification, accuracy, precision, recall, F1, ROC-AUC, PR-AUC and the lift in the top decile were reported.
For clustering, the silhouette coefficient, the Davies--Bouldin index and the Calinski--Harabasz index were
reported. Anomaly detectors were evaluated by precision and recall on a copy of the data containing planted
errors of known type.

\subsection{Benchmarks}
Following \citet{makridakis2018}, every forecasting model is compared with two naive rules: \emph{naive}
(next month = this month) and \emph{seasonal naive} (next month = the same month last year). At company level,
Holt--Winters and SARIMA serve as statistical benchmarks. The majority-class rule is the benchmark for
classification.

\section{Statistical analysis}
The exploratory phase used descriptive statistics, Spearman rank correlation (robust to the skewed
distributions found), Kruskal--Wallis tests for differences between groups, and chi-square tests of
independence. Because the sample is very large, almost any difference is statistically significant.
\emph{Effect sizes} ($\varepsilon^2$, $\rho^2$ and Cram\'er's $V$) were therefore reported and used to judge
practical importance. A fixed-effects log--log regression estimated the price elasticity of demand
\citep{wooldridge2020}.

\section{Two complementary modelling studies}
The predictive work was carried out twice, with different aims. The \emph{production study}
(\cref{ch:models}) builds the best possible models for the business: rich features, randomised hyper-parameter
search over expanding-window folds, statistical benchmarks and forecast reconciliation. The
\emph{algorithm-level study} (\cref{ch:algorithms}) examines each of the nine algorithms named in the proposal
separately, with a simpler shared feature set, a single 2024 validation year and one transparent script per
algorithm. Its aims are to explain how each algorithm learns and fails, and to check the production results
with an independent implementation.

\section{Tools and environment}
The framework was implemented in Python~3 with pandas and NumPy for data handling \citep{mckinney2022},
SciPy and statsmodels for statistics and time-series models, scikit-learn for machine learning
\citep{pedregosa2011}, XGBoost for gradient boosting \citep{chen2016}, Matplotlib and Plotly for
visualisation, SQLite for the dimensional data store, and Streamlit for the web dashboard. The code is
version-controlled with Git and published on GitHub. The dashboard is deployed on Streamlit Community Cloud and
the website on GitHub Pages. The report was typeset in \LaTeX. All random processes use a fixed seed (42), so every result can be reproduced exactly.

\section{Ethical considerations}
The dataset contains no names, contact details or other personally identifiable information. Customers are
described only by type (for example, \emph{Retailer}) and region. The data were used solely for academic
purposes, in line with the ethical principles stated in the proposal. Model limitations and uncertainties are
reported transparently (\cref{sec:limitations}), so that forecasts are not over-trusted in decision-making.
